\SourceLeaf{432}%
\[
\begin{aligned}
Z&=\pm\frac{R}{\BareRoot{(PP+QQ+RR)}}\\
&=\pm\frac{1}{\BareRoot{(1+tt+uu)}},
\end{aligned}
\]
also
\[
AC-BB=\frac{TV-UU}{(1+tt+uu)^{2}}.
\]

So wie \textls[100]{jedem} Punkt der krummen Fläche vermöge der daselbst gezogenen Normale und des dieser parallelen Radius der Hülfskugel ein besonderer Punkt $L$ auf der Fläche der Hülfskugel correspondirt, so wird aus den, allen Punkten einer \textls[100]{Linie} auf der krummen Fläche auf der Oberfläche der Hülfskugel correspondirenden, Punkten eine Linie gebildet, die jener correspondirt, und eben so wird jeder begrenzten Figur auf der krummen Fläche eine begrenzte Figur auf der Hülfskugel entsprechen; die Area auf der letztern soll als das Mass der Amplitudo der erstern betrachtet werden: wir werden diese Area entweder als eine Zahl betrachten, wobei das Quadrat des Halbmessers der Hülfskugel als Einheit zum Grunde liegt, oder auch sie durch Grade u.\,s.\,w. ausdrücken, indem wir die Area der Halbkugel $=360^{0}$ setzen.

Die Vergleichung der Area auf der krummen Fläche mit der entsprechenden Amplitudo führt auf den Begriff von dem, was wir das Krümmungsmass der Fläche nennen. Ist jene dieser proportional, so soll die Krümmung gleichförmig heissen, und der Quotient, wenn man die Amplitudo mit der Fläche dividirt, soll das Krümmungsmass heissen: diess ist der Fall, wenn die krumme Fläche eine Kugel ist, und das Krümmungsmass ist dann ein Bruch, dessen Zähler $1$, der Nenner das Quadrat des Halbmessers ist.

Wir werden das Krümmungsmass als positiv betrachten, wenn die Grenzen um die Figuren auf der krummen Fläche und auf der Hülfskugel in einerlei Sinn laufen, negativ, wenn die Grenzen im entgegengesetzten Sinn die Figuren umgeben. Findet jene Proportionalität nicht statt, so ist die Fläche ungleichförmig gekrümmt, und an jeder Stelle findet ein besonderes Krümmungsmass statt, welches aus der Vergleichung unendlich kleiner auf der krummen Fläche und der Hülfskugel zusammengehörender Theile hervorgeht. Es sei $d\sigma$ ein Flächenelement auf jener, $d\Sigma$ das correspondirende auf der Hülfskugel, so wird
\[
\frac{d\Sigma}{d\sigma}
\]
das Krümmungsmass an dieser Stelle sein.
